\section{A universal Euler contribution at every rational prime}
\label{sec:s5-universal-euler}

The preceding variational argument can retain a single
lower bound for all of the remaining Euler terms,
using the discriminant to distinguish ramified from
unramified primes. This is one uniform estimate from
the finite \(S_5\)-cover; it is not a sequence of
separate retained-prime estimates.
Keep the actual normal closure \(L/\mathbb Q\),
\(\Delta=\Delta_L\), \(H_\Delta\), \(F_{\Delta,k}\),
and \(J_{\Delta,k}\) of MA-RB-007.

\paragraph{MA-RB-008 (the universal all-prime variational bound).}
For \(\sigma>1\), define
\[
 \begin{split}
 \mathcal B_\Delta(\sigma)
   &=120\sum_{\substack{p\ {\rm prime}\\p\ne5}}
                   \frac{\log p}{p^{6\sigma}-1}
        -114\sum_{p\mid\Delta}\frac{\log p}{p^{6\sigma}-1}\\
   &=-120\,\frac{\zeta_{\mathbb Q}'}{\zeta_{\mathbb Q}}(6\sigma)
       -\frac{120\log5}{5^{6\sigma}-1}
        -114\sum_{p\mid\Delta}\frac{\log p}{p^{6\sigma}-1},\\
 \mathcal H_\Delta(\sigma)
   &=H_\Delta(\sigma)-\mathcal B_\Delta(\sigma).
 \end{split}
 \tag{MA-RB.85}
\]
The prime on \(\zeta_{\mathbb Q}\) means differentiation
with respect to its argument. The hypothesis at five
gives \(5\nmid\Delta\). Thus \(\mathcal B_\Delta\)
retains coefficient \(120\) at every prime not dividing
\(5\Delta\), and coefficient \(6\) at every prime
dividing \(\Delta\). The prime five is excluded, since
its exact full Euler contribution is already retained
in \(H_\Delta\).

For every actual field in the specified family,
\[
             0<\mathcal H_\Delta(\sigma)<H_\Delta(\sigma)
                                  \qquad(\sigma>1).
 \tag{MA-RB.86}
\]
For each integer \(k\ge1\), put
\[
 \begin{aligned}
 \mathcal F_{\Delta,k}(\sigma)
   &=(\sigma-\tfrac12)^2
           -\frac{9k(2\sigma-1)}{\mathcal H_\Delta(\sigma)},\\
 \mathcal J_{\Delta,k}
   &=\inf_{\sigma>1}\mathcal F_{\Delta,k}(\sigma).
 \end{aligned}
 \tag{MA-RB.87}
\]
This infimum is finite and attained in \((1,\infty)\), and
\[
                  \mathcal J_{\Delta,k}<J_{\Delta,k}<\tfrac14.
 \tag{MA-RB.88}
\]
If the comparison-failure branch of MA-RB-002 occurs
with its actual pole order \(k\) at \(\gamma\), then
\[
 \begin{gathered}
 9k\left(\frac1{\sigma-\gamma}
             +\frac1{\sigma-1+\gamma}\right)
                 \le\mathcal H_\Delta(\sigma)
                                  \qquad(\sigma>1),\\
                  (\gamma-\tfrac12)^2
                            \le\mathcal J_{\Delta,k}.
 \end{gathered}
 \tag{MA-RB.89}
\]
In particular \(\mathcal J_{\Delta,1}\le0\) excludes
that branch. Otherwise the current shared-real-zero
criterion can be confined to
\[
 \begin{gathered}
 B(\gamma)=0,\qquad
       L(\gamma,V\eta)L(\gamma,X)=0,\\
 \tfrac12<\gamma
       \le\tfrac12+\sqrt{\mathcal J_{\Delta,1}}<1 .
 \end{gathered}
 \tag{MA-RB.90}
\]
If the previous interval in MA-RB-007 remained nonempty,
this either excludes its branch or strictly decreases
its upper endpoint. Neither the numerical exclusion
criterion nor nonvanishing on the remaining interval
is asserted for every discriminant in the family.

\begin{proof}
\emph{Residue degrees at ramified as well as unramified primes.}
Fix a rational prime \(p\) and a prime \(\mathfrak p\)
of \(L\) above it. Let \(D_{\mathfrak p}\) and
\(I_{\mathfrak p}\) be its decomposition and inertia
groups in \(G=\operatorname{Gal}(L/\mathbb Q)=S_5\).
Their residue-field exact sequence is
\[
 1\longrightarrow I_{\mathfrak p}
  \longrightarrow D_{\mathfrak p}
  \longrightarrow
       \operatorname{Gal}(\mathbb F_{p^{f_p}}/\mathbb F_p)
  \longrightarrow1 .
 \tag{MA-RB.91}
\]
Here \(f_p\) is the residue degree. The last group
is generated by the arithmetic residue Frobenius
\(u\mapsto u^p\) and has order \(f_p\). Choose any
lift \(d_p\in D_{\mathfrak p}\) of that generator.
The order of its image divides the order of \(d_p\).
The possible element orders in \(S_5\), by the cycle
types \(1,2A,2B,3A,6A,4A,5A\), are respectively
\(1,2,2,3,6,4,5\). Therefore
\[
                     1\le f_p\le6
                         \quad\hbox{for every rational prime }p.
 \tag{MA-RB.92}
\]
No assumption that \(I_{\mathfrak p}\) is trivial
has entered this argument. In particular it applies
to primes dividing the discriminant.

Because \(L/\mathbb Q\) is Galois, all primes over
\(p\) have this residue degree, a common ramification
index \(e_p\), and a common decomposition-group order.
Their number is \(g_p=120/(e_pf_p)\ge1\).
The local factor of the Dedekind zeta function is
\((1-p^{-f_ps})^{-g_p}\); its contribution to the
positive logarithmic derivative on the real
half-line is
\[
 E_p(\sigma)
     =\frac{g_pf_p\log p}{p^{f_p\sigma}-1}.
 \tag{MA-RB.93}
\]
This includes every power of every prime ideal
over \(p\), also when it is ramified.

For a fixed \(x>1\), the function
\(u\mapsto u/(x^u-1)\) is strictly decreasing on
\(u>0\). Indeed, after multiplying its derivative
by the positive denominator, its sign is that of
\(e^v(1-v)-1\), where \(v=u\log x>0\).
This last expression is zero at \(v=0\) and has
derivative \(-ve^v<0\).
Use \(x=p^\sigma\), \(f_p\le6\), and \(g_p\ge1\)
in (MA-RB.93). It follows that
\(E_p(\sigma)\ge6\log p/(p^{6\sigma}-1)\) at every
prime. At an unramified prime there is the stronger
identity \(g_pf_p=120\). Since \(f_p\le6\),
\[
 E_p(\sigma)=\frac{120\log p}{p^{f_p\sigma}-1}
                  \ge\frac{120\log p}{p^{6\sigma}-1}.
\]
Combining the two cases gives the uniform bound
\[
 E_p(\sigma)\ge c_p\frac{\log p}{p^{6\sigma}-1},
 \qquad
 c_p=
 \begin{cases}
 120,&p\nmid\Delta,\\
 6,&p\mid\Delta,
 \end{cases}
 \qquad \sigma>1.
 \tag{MA-RB.94}
\]
This is a bound for the full local Euler contribution,
not a claim that a ramified Frobenius itself has a
specified conjugacy class in \(S_5\). The comparison
is pointwise on the real half-line. It is not
coefficientwise positivity of the difference of
Dirichlet series, nor an inclusion of Artin motives.

\emph{Strict positivity after the universal subtraction.}
All the positive Euler sums just used converge
absolutely on \(\sigma>1\); the comparison series
in (MA-RB.85) converges already when \(6\sigma>1\).
The ordinary rational zeta Euler product gives
the logarithmic-derivative expression there, with
exactly the stated coefficient \(120\), the
exclusion of five, and the finite ramified correction
of \(120-6=114\).
Subtracting \(\mathcal B_\Delta\) from the exact identity
(MA-RB.77) gives
\[
 \begin{split}
 \mathcal H_\Delta(\sigma)
   ={}&\sum_\rho\operatorname{Re}\frac1{\sigma-\rho}\\
      &+\sum_{p\ne5}
          \left(E_p(\sigma)
                  -c_p\frac{\log p}{p^{6\sigma}-1}\right).
 \end{split}
 \tag{MA-RB.95}
\]
The zero sum includes all nontrivial completed
zeros with multiplicity and is absolutely convergent
in this real-part form, as proved in MA-RB-001.
Its terms are strictly positive for each zero in
the usual containing strip
\(0\le\operatorname{Re}\rho\le1\).
Every Euler remainder is nonnegative by
(MA-RB.94).

We verify that the completed divisor is nonempty,
without invoking any particular zero or zero location.
If it were empty, the same exact Hadamard identity
of MA-RB.8--.12 would give, for every \(\sigma>1\),
\[
 0=\frac1\sigma+\frac1{\sigma-1}
       +\frac12\log\Delta
       +60\bigl(\psi(\sigma)-\log(2\pi)\bigr)
       +\frac{\zeta_L'}{\zeta_L}(\sigma).
\]
The Euler series implies
\(\zeta_L'/\zeta_L(\sigma)\to0\) as \(\sigma\to\infty\):
each positive term tends to zero, and for
\(\sigma\ge2\) the series at \(2\) is an absolutely
summable majorant. The Jensen estimate already proved
in MA-RB-007 gives
\(\psi(\sigma)=\log\sigma+O(1/\sigma)\).
The right side would therefore be
\(60\log\sigma+O_\Delta(1)\), a contradiction.
The completed divisor is consequently nonempty.
Its sum in (MA-RB.95) is strictly positive at every
\(\sigma>1\). This proves
\(\mathcal H_\Delta(\sigma)>0\) even if every retained
Euler remainder happens to vanish.

Also \(\mathcal B_\Delta(\sigma)>0\), since its
positive-coefficient form is
\[
 120\sum_{p\nmid5\Delta}\frac{\log p}{p^{6\sigma}-1}
       +6\sum_{p\mid\Delta}\frac{\log p}{p^{6\sigma}-1}.
\]
This proves the other strict inequality in (MA-RB.86).
Discarding the discriminant information would instead
retain only \(6\sum_{p\ne5}\log p/(p^{6\sigma}-1)\).
The present single bound is strictly larger, because
there is an unramified prime different from five.
Only the discriminant-aware bound is used in the
minimum (MA-RB.87).

\emph{The retained real pair and the minimum.}
Under the hypothetical comparison failure, the
distinct real zeros \(\gamma,1-\gamma\) each
have multiplicity at least \(9k\), as proved
in MA-RB-002/007. Retain them in (MA-RB.95)
and discard only nonnegative terms. This gives
the first line of (MA-RB.89).
With \(a=\gamma-1/2\) and \(x=\sigma-1/2\),
their contribution is \(18kx/(x^2-a^2)\).
All denominators are positive. The same algebra
as (MA-RB.78) gives
\(a^2\le\mathcal F_{\Delta,k}(\sigma)\), and
taking the infimum proves the second line.

The function \(\mathcal B_\Delta\) is smooth on a
neighborhood of \(\sigma=1\) and tends to zero
at infinity. Indeed its differentiated prime
series converge uniformly on every closed
subinterval of \(6\sigma>1\).
Consequently the exact endpoint estimates of
MA-RB-007 become
\[
 \begin{aligned}
 \mathcal H_\Delta(1+t)&=t^{-1}+O_\Delta(1),\\
 \mathcal F_{\Delta,k}(1+t)
       &=\tfrac14+(1-9k)t+O_{\Delta,k}(t^2),\\
 \mathcal H_\Delta(\sigma)
       &=60\log\sigma+O_\Delta(1),\qquad
 \mathcal F_{\Delta,k}(\sigma)\longrightarrow+\infty .
 \end{aligned}
 \tag{MA-RB.96}
\]
The first two assertions concern \(t\downarrow0\);
the last line concerns \(\sigma\to\infty\).
The digamma estimate needed there is already
proved by the two Jensen inequalities in
MA-RB-007, so no new gamma asymptotic is assumed.
Continuity and strict positivity of
\(\mathcal H_\Delta\), the negative linear
coefficient \(1-9k\), and (MA-RB.96) imply
that the infimum in (MA-RB.87) is finite
and attained on a compact subinterval of
\((1,\infty)\), with value less than \(1/4\).

To obtain the strict improvement, let
\(\sigma_*\in(1,\infty)\) attain the old
minimum \(J_{\Delta,k}\), whose existence
is part of MA-RB-007. The strict positive
denominator comparison (MA-RB.86) gives
\[
 \mathcal J_{\Delta,k}
   \le\mathcal F_{\Delta,k}(\sigma_*)
   <F_{\Delta,k}(\sigma_*)=J_{\Delta,k}.
\]
This proves (MA-RB.88), including strictness,
without a numerical optimization or a
compactness assumption about a limiting
minimizer.

Finally, \(\mathcal F_{\Delta,k}\) decreases
pointwise with \(k\), so
\(\mathcal J_{\Delta,k}\le\mathcal J_{\Delta,1}\).
If the latter is nonpositive, (MA-RB.89)
contradicts \(\gamma>1/2\). Otherwise
(MA-RB.89) gives the interval in (MA-RB.90),
and MA-RB-002 supplies its two vanishing
factors. This proves all stated conclusions.
\end{proof}

\paragraph{Finite-cover and Artin-motive meaning of the local bound.}
The representation underlying \(\zeta_L\) is the regular
permutation representation \(\mathbb Q[G]\), the
Artin motive of the finite cover \(\operatorname{Spec}L\)
over \(\operatorname{Spec}\mathbb Q\).
At a possibly ramified prime, take
\(\mathbb Q[G]^{I_{\mathfrak p}}\).
For the left regular action, an invariant basis is
indexed by the \(I_{\mathfrak p}\)-orbits \(I_{\mathfrak p}h\).
There are \(120/e_p=g_pf_p\) of them.
The chosen lift \(d_p\) normalizes inertia and
permutes these orbits. Its \(r\)th power fixes
one if and only if \(d_p^r\in I_{\mathfrak p}\),
which is equivalent to \(f_p\mid r\).
Every cycle therefore has length \(f_p\),
and there are exactly \(g_p\) cycles. Thus
\[
 \det(1-Td_p\mid\mathbb Q[G]^{I_{\mathfrak p}})
                         =(1-T^{f_p})^{g_p}.
\]
This proves directly the Artin local-factor
interpretation of (MA-RB.93).
We used arithmetic residue Frobenius. Replacing
it by its inverse gives the same determinant
for these complete cyclic permutation blocks.
The argument concerns the finite cover and
its local inertia invariants; it does not
assert an effective or positive realization
of an unknown analytic zero divisor.

\paragraph{Established locators and scope.}
Calegari's Theorem~1.1, assumptions (1)--(2),
supplies the actual \(S_5\) normal closure,
signature and exact prime-five datum; the
supplied primary LaTeX locator is
\texttt{S5ArtinWeb.tex}, lines 123--136.
The standard Galois decomposition-group
sequence in (MA-RB.91) fixes the residue-field
convention, and the local regular-representation
calculation above explains every ramified
Euler coefficient used here.

For the finite-cover interpretation, the exact
inherited locators are SGA~1, Expos\'e~V, \S7,
and Deligne (1979), \emph{Valeurs de fonctions
\(L\) et p\'eriodes d'int\'egrales},
Definition~6.1 and \S\S6.2--6.3; equation
(0.5.3) supplies the previously fixed
restriction-of-scalars/induction convention.
Those categorical locators are the established
ones in MA-S5; no new claim that they prove
the analytic inequality is made.
The exact Hadamard identity and constant are
MA-RB.8--.12, the shared-zero multiplicity is
MA-RB-002, and the preceding variational
construction is MA-RB-007. The new proof uses
no further functoriality, critical-line
location, \(\mathsf A\), \(M\), or
\((\mathcal R_B)\).
